\documentclass[10pt,a4paper]{amsart}
\usepackage[margin=19mm]{geometry}
\usepackage{amsmath,amssymb,mathtools}
\usepackage[scr=rsfs,cal=euler]{mathalfa}
\usepackage{xfrac}
\usepackage[unicode,hidelinks,bookmarks=false]{hyperref}
\newcommand{\ie}{i.e.\ }
\newcommand{\cf}{cf.\ }
\newcommand{\resp}{respectively\ }
\newcommand{\Subsection}{\subsection}
\providecommand{\nobreakdash}{\nobreak\hbox{-}}
\setlength{\emergencystretch}{2em}
\multlinegap0pt
\usepackage{bm}
\usepackage{bbm}
\usepackage{dsfont}
\usepackage{xspace}
\usepackage{cancel}
\usepackage{mathtools}
\usepackage{amsthm}
\makeatletter
\@ifundefined{lemma}{\newtheorem{lemma}{Lemma}}{}
\makeatother
\newcommand{\cC}{\mathcal{C}}
\newcommand{\cG}{\mathcal{G}}
\newcommand{\cH}{\mathcal{H}}
\newcommand{\cJ}{\mathcal{J}}
\newcommand{\cL}{\mathcal{L}}
\title[A general-position problem for planar line arrangements]{A general-position problem for planar line arrangements}
\author{Oliver Roche-Newton}
\date{Source: arXiv:2607.25742v1 (CC BY 4.0)}
\begin{document}
\maketitle

\noindent\emph{Source and licence.} A general-position problem for planar line arrangements. Version arXiv:2607.25742v1 (CC BY 4.0), distributed under the Creative Commons Attribution 4.0 International licence, as recorded in the arXiv licence metadata for this exact version and shown on the abstract page https://arxiv.org/abs/2607.25742v1. Full rights notice: licenses/arxiv-2607.25742-CC-BY-4.0.txt.\par\smallskip

\noindent\textit{Editorial scope.} Counted unit: the Lemma whose source label is \texttt{lem:cont2} in the article (source0.tex lines 372--385, source label \texttt{lem:cont2}), delivered together with the article's own complete original proof of it (source0.tex lines 387--428, one \texttt{proof} environment of 42 lines). No mathematics is written, repaired or re-derived: statement and proof are the article's own text line for line. Required context retained uncounted: Ldefn (lines 194--196); lem:sslL (lines 199--204); lem:cont1 (lines 323--336). Every definition, display and label of those lines is retained unchanged. Omitted from this package and not redistributed: 1--193, 197--198, 205--322, 337--371, 386--386, 429--652. Nothing inside a retained excerpt is omitted or altered: every retained body is delivered exactly as the article's lines are written, so no omission receipt of any kind accompanies this package. Citations: the retained excerpts carry no citation command at all, so no bibliography entry of the article is transcribed and no reference is redistributed. Reference adaptation: none was needed and none was made; every reference inside the delivered unit resolves inside it, so no unresolved reference or citation is produced. Printed slips kept literal: no printed word, symbol, hypothesis or number of the counted statement or of its proof was altered. Layout and class: the article's own class is \texttt{amsart} with packages including amsthm, amssymb, latexsym, amsmath, mathtools, color, geometry, xcolor, hyperref, tikz, tikz-cd, comment, caption, graphicx; the wrapper is the lane's pinned \texttt{amsart} editorial wrapper at 10pt on A4, which restates the theorem-like environments the retained text needs and adds the article's own macro definitions that the retained text needs (\texttt{\textbackslash cC}, \texttt{\textbackslash cG}, \texttt{\textbackslash cH}, \texttt{\textbackslash cJ}, \texttt{\textbackslash cL}), with extra wrapper packages (bm, bbm, dsfont, xspace, cancel, mathtools). The delivered numbering is the unit's own. Assets: no retained span contains a figure environment, a graphics-inclusion command, a PDF-page inclusion, a table or a TikZ picture; no image file is copied into this package, the package declares no asset dependency and no image file is redistributed with it. Licence: version arXiv:2607.25742v1 of this article is distributed under the Creative Commons Attribution 4.0 International licence, as recorded in the arXiv licence metadata for this exact version and shown on the abstract page https://arxiv.org/abs/2607.25742v1. A full rights notice is delivered at licenses/arxiv-2607.25742-CC-BY-4.0.txt. Characters: every retained excerpt is pure ASCII, so the delivered engine, XeLaTeX with its default Latin Modern text font, reports no missing character.
\begin{equation} \label{Ldefn}
\mathcal L:= \{ \ell_{a,b} : a \in [r], b \in [r^7] \},
\end{equation}

\begin{lemma} \label{lem:sslL}
Let $\mathcal J \subseteq \mathcal L$ such that $|\mathcal J| \geq 8r^7$. Then
\[
|T(\mathcal J)| \geq \frac{|\mathcal J|^3}{1000r^8}.
\]
\end{lemma}

\begin{lemma} \label{lem:cont1}
   Let $r$ be a sufficiently large integer and let $\cL$ be the line set defined in \eqref{Ldefn}. Let $\mathcal J \subseteq \mathcal L$ such that $|\cJ| \geq 8r^7$ and let $\cG=\mathcal H[ \mathcal J]$ be the subhypergraph of $\mathcal H$ induced by $\mathcal J$. For all $0<\delta <1$, there exists a family of sets $\mathcal C_{\mathcal J} \subseteq 2^{\mathcal J}$ such that the following three statements hold.
    \begin{itemize}
        \item If $A \subseteq \mathcal J$ is an independent set in $\cG$ then there exists $C \in \mathcal C_{ \mathcal J}$ such that $A \subseteq C$.
        \item For all $C \in \mathcal C_{ \mathcal J}$,
        \[
        |E(\cH[C])| \leq r^{-\delta} |E( \cH[\cJ])|
        \]
        \item For some absolute constant $c'$, we have 
        \[
        \log | \cC_{\cJ}| \leq c'\delta r^{4+\delta} \log^2 r .
        \]
    \end{itemize}
\end{lemma}

\begin{lemma} \label{lem:cont2}
   Let $0 < \delta < 1$ and let $r$ be sufficiently large. Let $\cL$ be the set of lines defined in \eqref{Ldefn} and let $\cH$ be the $3$-uniform hypergraph encoding intersecting triples of lines from $\cL$. There exists a set $\cC \subseteq 2^{\cL}$ such that the following three statements hold.
   \begin{itemize}
        \item If $A \subseteq \mathcal L$ is an independent set in $\cH$ then there exists $C \in \mathcal C$ such that $A \subseteq C$.
        \item For all $C \in \mathcal C$,
        \[
        |C| \leq 8r^7.
        \]
        \item For some constant $c$, we have 
        \[
        \log | \cC| \leq c\delta r^{4+\delta} \log^2 r .
        \]
    \end{itemize}
\end{lemma}

\begin{proof}

%\orn{Add a picture to this proof showing the tree structure.}


    Iteratively apply Lemma \ref{lem:cont1}. We begin by applying the lemma with $\cJ=\cL$ to obtain a set of containers $\cC_{1}$ satisfying 
    \[
    \log | \cC_1| \leq c' \delta r^{4+\delta} \log^2 r
    \]
    such that any independent set $A$ in $\cH$ satisfies $A \subseteq C$ for some $C \in \cC_{1}$. For each $C \in \cC_{1}$, there are two possibilities.
    \begin{itemize}
    \item If $|C| > 8r^7$ we apply Lemma \ref{lem:cont1} with $\cJ=C$. We obtain a family $\cC_{C}$ of subsets of $C$ such that every independent set in $\cH[C]$ is contained inside a set in $\cC_{C}$.
        \item If $|C| \leq 8r^7$ then we put $C$ into our final set of containers $\cC$. For the purpose of describing this algorithm, it is convenient to define $\cC_{C}= \{ C \}$.
    \end{itemize}
We define 
\[
\cC_2:= \bigcup_{C \in \cC_1} \cC_C
\]
Note that $\mathcal C_2$ is a container set for $\cH$. Indeed, suppose that $X$ is an independent set in $\cH$. Then there is some $C \in \mathcal C_1$ such that $X \subseteq C$. Also, $X$ is an independent set in the hypergraph $\cH[C]$, which implies that $X \subseteq C'$ for some $C' \in \mathcal C_C \subseteq \mathcal C_2$.
    
    We then repeat this process for each set in $\cC_2$ to obtain a larger family of containers $\cC_3$ with smaller component sets. We can visualise this as a tree of containers, where the final set of containers $\cC$ consists of all of the leaf nodes. This process terminates when all of the leaf nodes correspond to sets $C$ with $|C| \leq 8r^7$. 
    
    We claim that this process terminates after at most $\lceil \frac{4}{\delta} \rceil$ steps. Indeed, let $C$ be a node in this tree at depth at least $\lceil \frac{4}{\delta}\rceil $. Since each application of Lemma \ref{lem:cont1} reduces the number of edges by a factor of $r^{-\delta}$, it follows that the number of edges determined by $C$ satisfies
    \begin{equation} \label{edges}
    |E(\cH[C])| \leq c_1r^{16} \log r(r^{-\delta})^{\lceil \frac{4}{\delta}\rceil }\leq c_1r^{12} \log r.
    \end{equation}
     Now suppose for a contradiction that $|C| > 8r^7$. Lemma \ref{lem:sslL} then implies that
    \[
   | E(\cH[C])| \geq \frac{|C|^3}{1000r^8}
    \]
    and combining this with \eqref{edges} and rearranging yields
    \[
    |C| \leq c_3 r^{20/3} \log^{1/3} r
    \]
    for an absolute constant $c_3$. But $|C| > 8r^7$ and this is a contradiction for $r$ sufficiently large.

    At the end of this process, we have a set of containers $\cC$ such that the bound $|C| \leq 8r^7$ holds for all $C \in \cC$. Finally, it follows from the upper bound for the size of the container families given in Lemma \ref{lem:cont1} that
    \[
    |\cC| \leq \left(2^{c'\delta r^{4+\delta} \log^2 r} \right)^{\lceil \frac{4}{\delta}\rceil } \leq  2^{cr^{4+\delta} \log^2 r},
    \]
as required.
\end{proof}

\end{document}
